\begin{textsample}{NEG Dim 4 – 2010\_c6 – Score -2.00 – 2010\_c6/enade\_2010\_fisio\_fd971bea85.txt}  \label{ex:f4_neg_045}
CAPÍTULO 3 ANÁLISE TÉCNICA DA PROVA Este capítulo tem por objetivo apresentar o desempenho dos estudantes de Fisioterapia no ENADE/2010.
Para isso, foram calculadas as estatísticas básicas da prova em seu todo, bem como as estatísticas dos componentes relacionadas à Formação Geral e ao de Conhecimento Específico da Área.
Nas tabelas, são apresentados o tamanho da população inscrita e de presentes, e as seguintes estatísticas das notas⁷ : média do desempenho na prova, erro padrão da média, desvio padrão, nota mínima, mediana, nota máxima, coeficiente de assimetria e de variação.
As estatísticas apresentadas neste capítulo contemplam o total de estudantes e, separadamente, os \textbf{Ingressantes} e os Concluintes.
Foram calculadas tendo -se em vista as seguintes agregações : ( a ) as Grandes Regiões e o país como um todo ; ( b ) a Categoria Administrativa ; e ( c ) a Organização Acadêmica.
Em relação aos gráficos de distribuição de notas, o intervalo considerado foi de 10 unidades, aberto à esquerda e fechado à direita, com exceção do primeiro intervalo, [ 0 ; 10 ], fechado em ambos os extremos.
Para os gráficos de distribuição das notas das questões discursivas, foram consideradas mais duas categorias : questão em branco e nota zero. 3.1 ESTATÍSTICAS BÁSICAS DA PROVA 3.1.1 Estatísticas Básicas Gerais A Tabela 3.1 apresenta as estatísticas básicas da prova por grupo de estudantes.
A população total de estudantes regularmente inscritos⁸ para a prova foi 38.303 Destes, 31.024 estiveram presentes, sendo 19,0 \% o índice de não comparecimento.
A abstenção foi maior entre \textbf{Ingressantes} ( 27,1 \% ) do que entre Concluintes ( 4,8 \% ).
A nota média de toda a prova foi 36,4, sendo que os \textbf{Ingressantes} obtiveram média mais baixa, 30,7, que os Concluintes, cuja média foi 44,0.
O desvio padrão foi 13,2, sendo o do grupo de \textbf{Ingressantes} menor ( 10,2 ) do que o do grupo dos Concluintes ( 12,8 ) ; os coeficientes de variação ( 33,2 \% no grupo dos \textbf{Ingressantes} e 29,1 \% no dos Concluintes ) indicam menor dispersão relativa nas notas de Concluintes.
A nota máxima foi 86,4, obtida por Concluintes, ao passo que, a maior nota obtida por \textbf{Ingressantes} foi 76,7. ⁷ Algumas dessas estatísticas estão definidas no Capítulo 1. ⁸ As estatísticas das notas apresentadas consideram apenas os estudantes regularmente inscritos e com resultados validados por ausência de erro administrativo.
Uma comparação entre o desempenho de \textbf{Ingressantes} e Concluintes também pode ser realizada a partir dos dados apresentados no Gráfico 3.1, com a distribuição de notas finais.
Destaca -se que as proporções mais elevadas de notas, considerando -se o intervalo posterior à nota 40,0, foram atingidas principalmente pelos Concluintes, cuja concentração das notas está à direita dos \textbf{Ingressantes}.
Já nas faixas de notas mais baixas, inverteu -se o quadro identificado nas faixas de notas mais altas.
Essa característica é corroborada pelo valor da nota mediana 30,3 no grupo dos \textbf{Ingressantes}, abaixo da encontrada no grupo dos Concluintes, 43,8, transparecendo uma superioridade de desempenho dos Concluintes em relação aos \textbf{Ingressantes} na prova em tela ( Tabela 3.1 ).
Os coeficientes de assimetria indicam comportamento simétrico para os Concluintes, e positivo ( 0,2 ) para \textbf{Ingressantes}.
O Gráfico 3.2 apresenta informações referentes à nota final de \textbf{Ingressantes} e Concluintes, comparando -se os resultados de acordo com as Grandes Regiões do país, a Categoria Administrativa e a Organização Acadêmica.
Considerando -se as notas médias dos alunos \textbf{Ingressantes} em cada Grande Região, observa -se que, em relação às notas aferidas em âmbito nacional ( média de 30,7 para \textbf{Ingressantes} e de 44,0 para Concluintes ), existe diferença estatisticamente significativa ao nível de 95 \% entre a menor média, obtida na região Norte ( 29,4 ) e a maior, na região Centro-Oeste ( 31,5 ).
Entre os Concluintes, também existe diferença estatisticamente significativa das médias ao nível de 95 \%.
A maior média de notas foi obtida pela região Sul ( 46,1 ), ao passo que, a menor foi identificada na região Norte ( 40,4 ).
As demais médias de notas foram 43,8 para as regiões Nordeste e Centro-Oeste ( ambas com o mesmo valor ) e para a região Sudeste ( 43,9 ).
Levando -se em conta os grupos dos estudantes em Categorias Administrativas, observa -se que a pontuação média mais elevada entre os Concluintes foi encontrada em instituições públicas ( 52,5 ), com desempenho acima da média nacional.
Já os Concluintes provenientes de instituições privadas obtiveram média 43,1, abaixo da média nacional.
Quanto aos \textbf{Ingressantes}, as notas mais elevadas, também acima da média nacional para esse grupo de estudantes, ocorreram nas instituições públicas ( 37,8 ).
Os \textbf{Ingressantes} vinculados a instituições privadas obtiveram nota média 29,7, abaixo da média nacional.
Tendo como foco as notas médias dos estudantes agrupados por Organização Acadêmica, distribuídas entre Universidades, Centros Universitários e Faculdades, para os Concluintes vinculados a Universidades, a média 46,4 ficou acima da média nacional.
Pode -se verificar a mesma realidade entre os \textbf{Ingressantes} de Universidades, cuja nota média foi 32,9, também acima da nacional para esse grupo.
Considerando -se as demais Organizações Acadêmicas, observa -se que nos Centros Universitários e Faculdades, o desempenho de \textbf{Ingressantes} e de Concluintes ficou abaixo da média nacional, com notas médias 42,9 e 42,1, respectivamente. \textbf{Ingressantes} das mesmas Organizações Acadêmicas obtiveram médias 29,7 e 29,2, respectivamente. 3.1.2 Estatísticas Básicas no Componente de Formação Geral A Tabela 3.2 apresenta as estatísticas básicas em relação ao componente da prova que avalia a Formação Geral dos estudantes.
Os Concluintes obtiveram desempenho médio 47,6, que foi superior ao dos \textbf{Ingressantes}, 42,9.
Quanto à variabilidade das notas nos dois grupos, o desvio padrão das notas obtidas foi 18,4 para \textbf{Ingressantes} e 17,9 para Concluintes.
Os dois grupos alcançaram notas máximas semelhantes, sendo 96,0 a maior nota obtida por pelo menos um \textbf{Ingressante}, e 97,0 a maior nota entre os Concluintes.
O Gráfico 3.3 propicia a comparação entre Concluintes e \textbf{Ingressantes} relativamente ao desempenho no componente de Formação Geral.
Nesse componente, tal como na prova como um todo, os Concluintes apresentaram desempenho superior.
Dos \textbf{Ingressantes}, 19,9 \% obtiveram notas na faixa entre 40 e 50 pontos, sendo essa a faixa modal.
Nas faixas de notas seguintes, mais elevadas, verifica -se queda gradual de ocorrência de estudantes \textbf{Ingressantes}.
Entre os Concluintes, houve percentual maior de estudantes no intervalo entre 50 e 60 pontos, perto dos 20,9 \%.
Também foi observada queda gradual no desempenho, nos intervalos de notas mais elevados, com cerca de 7,5 \% desses estudantes tendo atingido a faixa de 70 a 80 e 2,0 \%, a de notas entre 80 e 90 pontos.
Apenas 3,1 \% dos Concluintes tiveram desempenho na menor faixa de notas, que vai de zero a 10 pontos.
O coeficiente de assimetria foi negativo ( -0,2 ) para Concluintes, havendo comportamento simétrico para \textbf{Ingressantes}.
No Gráfico 3.4 são apresentadas as informações referentes ao desempenho dos \textbf{Ingressantes} e dos Concluintes no componente de Formação Geral, em diferentes agregações como Grande Região do país, Categoria Administrativa e Organização Acadêmica.
Considerando -se as notas médias dos estudantes segundo cada Grande Região, observa -se que no grupo de \textbf{Ingressantes} existe diferença estatisticamente significativa ao nível de 95 \% entre as médias nas regiões Centro-Oeste ( 44,2, a mais alta ) e Norte ( 41,7, a mais baixa ).
No grupo de Concluintes, a nota mais elevada ocorreu na região Sul, com média 49,2, enquanto a menor foi 46,3, correspondente à região Norte, diferença esta também estatisticamente significativa ao nível de 95 \%.
Essa região, por conseguinte, foi a que obteve as menores médias de notas nos dois agrupamentos de participantes.
Tendo como foco as Categorias Administrativas, as instituições públicas obtiveram médias mais elevadas do que as privadas : para Concluintes, notas médias de 54,6 em cursos de IES públicas e de 46,9 nos das privadas.
Para os \textbf{Ingressantes}, as notas foram, respectivamente, 54,4 e de 41,2.
Em comparação à média nacional, as notas de \textbf{Ingressantes} e de Concluintes vinculados às instituições públicas foram superiores, ao passo que, nas instituições privadas, foram inferiores.
Considerando -se o tipo de Organização Acadêmica, nota -se diferença estatisticamente significativa a 95 \% entre a menor e a maior média.
Nas Universidades, houve desempenho melhor do que em Centros Universitários e Faculdades.
Neste Neste primeiro tipo de Organização Acadêmica, \textbf{Ingressantes} tiveram nota média 46,2, enquanto Concluintes obtiveram média 49,8.
Nos Centros Universitários e Faculdades, os valores foram mais baixos : os \textbf{Ingressantes} tiveram média 41,2 nos Centros Universitários e 40,5 nas Faculdades ; entre Concluintes, a média foi 45,9 para Centros Universitários e 46,2 para as Faculdades. 3.1.3 Estatísticas Básicas do Componente de Conhecimento Específico A Tabela 3.3 apresenta as estatísticas básicas referentes à parte do componente de Conhecimento Específico da Área de Fisioterapia.
A média do desempenho foi 33,5.
Os Concluintes apresentaram desempenho médio ( 42,8 ) superior ao dos \textbf{Ingressantes} ( 26,7 ) ; as medidas de dispersão – desvio padrão e amplitude geral das notas – são menores no grupo dos \textbf{Ingressantes} do que no dos Concluintes.
Os coeficientes de variação, 38,6 \% no grupo dos \textbf{Ingressantes} e 32,0 \% no dos Concluintes, indicam dispersão relativa menor nas notas desses últimos.
As notas máximas obtidas por Concluintes e \textbf{Ingressantes} foram iguais a 77,3 e 88,7, respectivamente, notas inferiores àquelas encontradas no componente de Formação Geral ( 97,0 e 96,0, respectivamente ).
As medianas das distribuições das notas dos \textbf{Ingressantes} e dos Concluintes foram 27,0 e 42,5, respectivamente, enfatizando a superioridade dos Concluintes neste componente da prova.
Assim como os Gráficos 3.1 e 3.3, o Gráfico 3.5, apresentado a seguir, proporciona comparação do desempenho entre Concluintes e \textbf{Ingressantes}, em relação ao componente de Conhecimento Específico.
Dentre as três comparações feitas, observa -se que, embora a distribuição das notas seja semelhante à apresentada na prova como um todo, neste componente, a diferença de desempenho entre Concluintes e \textbf{Ingressantes} é mais acentuada, destacando -se o desempenho mais elevado dos primeiros.
Claramente a distribuição das notas dos \textbf{Ingressantes}, mais concentradas à esquerda em relação aos Concluintes, evidencia o percentual mais alto de notas nas faixas mais baixas para este grupo.
A maior concentração de Concluintes está no intervalo compreendido entre ( 30 e 40 ] pontos, ao passo que a de \textbf{Ingressantes} está entre ( 20 e 30 ] pontos.
Assim como os Gráficos 3.2 e 3.4, o Gráfico 3.6 apresenta uma comparação dos resultados em relação à Grande Região do país, à Categoria Administrativa e à Organização Acadêmica, agora levando em conta o desempenho dos \textbf{Ingressantes} e dos Concluintes no componente de Conhecimento Específico da prova.
Ao se observar as notas médias dos estudantes, segundo cada Grande Região do país, verifica -se que, entre os Concluintes, a mais elevada foi encontrada na região Sul ( 45,0 ), maior que a média nacional ( 42,8 ).
Em relação aos \textbf{Ingressantes}, também a região Sul apresenta a média mais elevada ( 27,6 ), seguida da região Centro-Oeste, com média igual a 27,3.
As médias mais baixas entre os Concluintes ( 38,4 ) e entre os \textbf{Ingressantes} ( 25,3 ) foram verificadas na região Norte.
Quanto à Categoria Administrativa, observa -se comportamento semelhante àquele da parte relativa à Formação Geral e à prova como um todo, ou seja, as médias mais elevadas tanto entre os Concluintes quanto entre os \textbf{Ingressantes} foram encontradas em instituições de ensino públicas.
Nesse tipo de IES, a nota média de Concluintes foi 51,8, enquanto a de \textbf{Ingressantes}, 32,2.
Tais médias estão situadas acima da média nacional.
Nas instituições privadas, as médias foram, mais uma vez, inferiores em relação às públicas.
Concluintes obtiveram média 41,9 e \textbf{Ingressantes}, 25,9.
Em relação à Organização Acadêmica, o desempenho dos Concluintes foi superior ao dos \textbf{Ingressantes} nas três categorias.
Nas Universidades, os Concluintes obtiveram a média mais elevada, igual a 45,3.
Abaixo estiveram os Centros Universitários ( 41,9 ) e as Faculdades ( 40,7 ).
Os \textbf{Ingressantes} de Universidades também tiveram nota um pouco mais elevada ( 28,4 ), seguidos pelos estudantes de Centros Universitários ( 25,8 ), que obtiveram média bem próxima aos de Faculdades ( 25,4 ).
Faculdades ( 40,7 ).
Os \textbf{Ingressantes} de Universidades também tiveram nota um pouco mais elevada ( 28,4 ), seguidos pelos estudantes de Centros Universitários ( 25,8 ), que obtiveram média bem próxima aos de Faculdades ( 25,4 ). 3.2 ANÁLISE DAS QUESTÕES OBJETIVAS 3.2.1 Componente de Formação Geral A Tabela 3.4 apresenta as estatísticas básicas relativas às oito questões objetivas do componente da prova que abrange a Formação Geral dos estudantes.
Como pode ser observado, os estudantes Concluintes obtiveram desempenho superior, com nota média 51,7, em relação aos alunos \textbf{Ingressantes}, cuja nota média foi 45,9.
As medidas de dispersão nos dois grupos de alunos foram similares, como indicam os desvios padrões, assim como a amplitude total das notas, sendo que a nota mínima foi zero e a máxima, 100,0, para os dois grupos de alunos avaliados.
Os coeficientes de variação, 46,2 \% no grupo dos \textbf{Ingressantes} e 39,7 \% no dos Concluintes, indicam menor dispersão relativa nas notas desses últimos.
A mediana foi igual a 57,1 no grupo dos Concluintes e 42,9 no dos \textbf{Ingressantes}.
O coeficiente de assimetria do total de estudantes revelou assimetria negativa ( -0,2 ).
A Tabela 3.5 apresenta as classificações das questões objetivas do componente de Formação Geral, segundo o índice de facilidade.
As questões objetivas da prova foram assim avaliadas : das oito questões, nenhuma teve o índice de facilidade classificado como muito fácil.
Duas questões foram tidas como fáceis, por terem índice de acertos situado na faixa entre 0,61 e 0,85 ( de 61,0 \% a 85,0 \% de acertos ).
Três questões foram consideradas de dificuldade média, situando -se no intervalo entre 0,41 e 0,60 do índice de facilidade, ou seja, houve entre 41,0 \% e 60,0 \% de acertos.
Uma questão foi classificada como difícil, situando -se no intervalo entre 0,16 e 0,40.
Por fim, duas questões apresentaram menos de 0,15 de acertos, razão pela qual foram classificadas como muito difíceis.
A Tabela 3.6 descreve os resultados das análises das questões objetivas relativas à Formação Geral, segundo o poder de discriminação, utilizando -se para tal, do índice de discriminação ponto bisserial.
Cinco das oito questões apresentaram índices acima de 0,40 e, assim, foram classificadas como muito boas nesta característica para esse grupo de alunos ; uma questão teve bom índice de discriminação, entre 0,30 e 0,39.
Foi classificada como de média discriminação uma das questões.
A questão de número 6 teve fraco nível de discriminação e foi eliminada do cômputo da nota final.
Na sequência, na Tabela 3.7, estão especificadas, por questão, os resultados da análise relativa aos índices de facilidade e de discriminação de cada questão no componente relativo à Formação Geral da prova de Fisioterapia.
O índice de facilidade variou de 0,07 a 0,73, e o de discriminação, de 0,19 a 0,53.
As questões com índices de discriminação muito bom, de números 1, 3, 4, 5 e 8, figuraram entre as fáceis ( questões 3 e 8 ) ou de dificuldade média ( questões 1, 4 e 5 ) desse conjunto.
A questão 8 foi a que apresentou maior poder discriminatório ( juntamente com a questão 3 ), com índice 0,53, e foi também a mais fácil, com uma proporção de 0,73 de acertos.
A questão de número 6 apresentou índice de facilidade 0,07, ou seja, um quantitativo pequeno de estudantes conseguiu resolvê -la, dentro do universo de participantes.
Esta questão obteve índice de discriminação 0,19 e, como já comentado, foi eliminada do cômputo da nota final.
A questão 7 obteve índice de discriminação médio, com 0,27, e a de número 2 obteve índice 0,35.
A Tabela 3.6 descreve os resultados das análises das questões objetivas relativas à Formação Geral, segundo o poder de discriminação, utilizando -se para tal, do índice de discriminação ponto bisserial.
Cinco das oito questões apresentaram índices acima de 0,40 e, assim, foram classificadas como muito boas nesta característica para esse grupo de alunos ; uma questão teve bom índice de discriminação, entre 0,30 e 0,39.
Foi classificada como de média discriminação uma das questões.
A questão de número 6 teve fraco nível de discriminação e foi eliminada do cômputo da nota final.
Na sequência, na Tabela 3.7, estão especificadas, por questão, os resultados da análise relativa aos índices de facilidade e de discriminação de cada questão no componente relativo à Formação Geral da prova de Fisioterapia.
O índice de facilidade variou de 0,07 a 0,73, e o de discriminação, de 0,19 a 0,53.
As questões com índices de discriminação muito bom, de números 1, 3, 4, 5 e 8, figuraram entre as fáceis ( questões 3 e 8 ) ou de dificuldade média ( questões 1, 4 e 5 ) desse conjunto.
A questão 8 foi a que apresentou maior poder discriminatório ( juntamente com a questão 3 ), com índice 0,53, e foi também a mais fácil, com uma proporção de 0,73 de acertos.
A questão de número 6 apresentou índice de facilidade 0,07, ou seja, um quantitativo pequeno de estudantes conseguiu resolvê -la, dentro do universo de participantes.
Esta questão obteve índice de discriminação 0,19 e, como já comentado, foi eliminada do cômputo da nota final.
A questão 7 obteve índice de discriminação médio, com 0,27, e a de número 2 obteve índice 0,35.
O Gráfico 3.7 analisa o comportamento da questão 8 de Formação Geral.
Trata -se da questão mais fácil e com o maior índice de discriminação desse componente da prova.
No gráfico a seguir, cada uma das cinco curvas representa o percentual de respostas em determinada alternativa da questão, em função da nota dos estudantes nesta parte da prova ( Formação Geral/Múltipla Escolha ) antes de possíveis eliminações pelo critério do ponto bisserial.
A curva em rosa corresponde à alternativa B, a correta para este item.
Assim, observa -se que entre os estudantes com notas mais baixas, desempenho mais fraco neste componente do exame, a situação mais frequente foi a escolha de uma das alternativas incorretas : a alternativa C ( em verde ) ou E ( em roxo ).
Na medida em que a nota aumenta, indicando desempenho melhor nesta parte da prova, aumenta concomitantemente a proporção de estudantes que selecionaram a alternativa correta B, atingindo 100 \% para as notas acima de 80,0.
Essa análise permite verificar como a questão discriminou entre os grupos de desempenho, justificando o alto índice obtido na questão.
Os gráficos relativos às demais questões de Formação Geral constam do Anexo I. 3.2.2 Componente de Conhecimento Específico A Tabela 3.8 apresenta as estatísticas básicas em relação às questões objetivas do componente de Conhecimento Específico da prova.
Nesse conjunto de questões, os estudantes Concluintes ( 47,3 ) obtiveram nota média superior aos \textbf{Ingressantes} ( 30,2 ).
Quanto à variabilidade das notas, o desvio padrão no grupo dos \textbf{Ingressantes} foi 11,7, e no grupo dos Concluintes, 14,8.
Os coeficientes de variação foram 38,7 \% e 31,3 \%, respectivamente, indicando menor dispersão relativa nas notas dos Concluintes.
Com relação à mediana, no grupo dos \textbf{Ingressantes}, esta foi igual a 31,8, estando pouco acima da média desse grupo, enquanto no grupo dos Concluintes, a mediana foi igual a 45,5, abaixo da média do grupo.
O coeficiente de assimetria foi 0,2 para \textbf{Ingressantes} e 0,1 para Concluintes.
A Tabela 3.9 apresenta as classificações das questões objetivas da parte da prova relativa ao componente de Conhecimento Específico, segundo o índice de facilidade.
Deve -se considerar o fato de as questões de números 13 e 29 terem sido anuladas pela Comissão, de modo que a classificação foi estabelecida com base em 25 das 27 questões.
A partir dos índices obtidos, pode -se concluir que a maioria das questões objetivas da prova foram tidas como difíceis : das 25 questões, quinze foram classificadas como difíceis e duas como muito difíceis.
Não houve questão classificada como muito fácil, ao passo que uma foi tida como fácil, na faixa de 0,61 a 0,85 do índice de facilidade, e sete consideradas médias, entre 0,41 e 0,60.
Na Tabela 3.10 são apresentados os índices de discriminação das questões objetivas do componente de Conhecimento Específico da prova.
Observa -se, segundo os índices estabelecidos, que a capacidade de discriminação das questões objetivas desse componente da prova foi classificada, para quatorze das 25 questões válidas, como boa, enquanto três delas tiveram índice de discriminação muito bom.
A maioria das questões – 17 em 25 – proporcionou, portanto, índices de discriminação bons ou muito bons.
Cinco delas foram classificadas como médias e três como fracas, sendo oito, por conseguinte, a quantidade de questões nos dois patamares mais baixos de discriminação.
Constata -se, assim, que a prova – no que se refere ao componente de Conhecimento Específico – possuía boa capacidade de discriminar entre aqueles que dominam ou não o conteúdo.
As questões 12, 16 e 33, com índice fraco de discriminação, foram eliminadas do cômputo das notas finais.
Na Tabela 3.11 estão especificados os resultados da análise relativa aos índices de facilidade e de discriminação de cada uma das questões do componente de Conhecimento Específico da prova de Fisioterapia.
As questões que alcançaram os maiores índices de discriminação, muito bom, foram as de números 20, 21, 30, classificadas como difíceis com respeito ao índice de facilidade.
Em termos de proporção, tais questões variaram de 0,30 a 0,47 quanto aos acertos.
Seus índices de discriminação situaram -se de 0,41 a 0,47.
A questão de número 12 foi a mais difícil dentre as 25 questões específicas válidas, com o mais baixo índice de facilidade, 0,11, o que, em termos percentuais, corresponde a 11,0 \% de acertos.
Essa questão apresentou poder discriminatório igualmente baixo, 0,06, o que comprova ter sido esta a mais difícil para os estudantes e foi eliminada do cômputo da nota final, como já comentado.
Destaca -se, também, a questão 15, cujo índice de facilidade foi 0,12, o que, em termos percentuais, corresponde a 12,0 \% de estudantes que responderam acertadamente.
Esta porém teve um índice de discriminação de 0,21.
O Gráfico 3.8 analisa o comportamento da questão 30 do componente de Conhecimento Específico.
Esta foi uma das questões mais fáceis da prova, apresentando índice de facilidade 0,47, ou seja, 47,0 \% dos estudantes assinalaram acertadamente a opção E, correspondente ao gabarito.
Seu índice de discriminação foi igual a 0,45, classificado como muito bom, embora outras questões tenham tido índices de discriminação mais elevados, como o caso da de número 20, com 0,47 de índice de discriminação.
Neste gráfico, cada uma das cinco curvas representa o percentual de respostas em determinada alternativa da questão 30, em função da nota dos estudantes nesta parte da prova antes de possíveis eliminações de questões pelo critério do ponto bisserial.
A alternativa correta E, representada no gráfico pela curva em roxo, foi escolhida em maiores proporções pelos alunos com desempenho melhor nesta parte da prova, enquanto as alternativas incorretas, também denominadas distratores, foram selecionadas principalmente por aqueles com notas mais baixas, ainda que aqueles com nota zero, na sua quase totalidade tenham deixado em branco ou marcaram mais de uma alternativa, comportamento considerado inválido.
A proporção de alunos que selecionaram a resposta correta E aumenta gradativamente com a nota dos alunos chegando a atingir 100 \% para notas acima de 80,0, enquanto a proporção dos que escolhem alternativas incorretas, decai a partir da primeira nota não nula como função da nota nesta parte da prova.
Os gráficos relativos às demais questões do componente de Conhecimento Específico constam do Anexo I. 3.3 ANÁLISE DAS QUESTÕES DISCURSIVAS 3.3.1 Componente de Formação Geral A análise dos resultados de desempenho dos estudantes de Fisioterapia, nas duas questões discursivas relativas à Formação Geral, encontra -se na Tabela 3.12 e no Gráfico 3.9.
Observa -se, a partir destes, que as notas médias foram mais baixas nesse conjunto do que no conjunto das objetivas.
A média geral dos \textbf{Ingressantes} nas questões objetivas de Formação Geral foi 45,9 ( Tabela 3.4 ), ao passo que nas questões discursivas, a média foi 38,3.
O mesmo ocorreu entre os Concluintes, que tiveram, em Formação Geral, média 51,7 nas questões objetivas e 41,4 nas questões discursivas.
Cabe destacar que, nas questões discursivas, o desempenho dos Concluintes foi pouco superior ao dos \textbf{Ingressantes}.
Os dois agrupamentos atingiram a mesma nota máxima, 100,0, o que acompanha a constatação de que não foi muito evidente uma superioridade dos Concluintes, em relação aos \textbf{Ingressantes}, nesta parte do exame.
Os coeficientes de assimetria indicam assimetria negativa para \textbf{Ingressantes} e Concluintes.
Verifica -se a semelhança das distribuições das notas desses grupos, no Gráfico 3.9, apresentado a seguir.
Na sequência, os resultados verificados para cada uma das questões discursivas de Formação Geral serão apresentados, estabelecendo -se relações com os conteúdos englobados em cada uma delas.
Os comentários da banca corretora a respeito do observado nas respostas dos estudantes, suas impressões e conclusões serão apresentados junto à análise de cada questão.
Cumpre esclarecer que, tendo em vista que as questões discursivas de Formação Geral são padronizadas, ou seja, constam de todas as provas, os comentários da banca são os mesmos para todas as carreiras acadêmicas, sendo direcionados a todos os estudantes que participaram do ENADE/2010.
A seguir, serão analisados os desempenhos de \textbf{Ingressantes} e Concluintes da Área de Fisioterapia nas duas questões discursivas de Formação Geral do ENADE/2010. 3.3.1.1 Análise da Questão Discursiva 9 do Componente de Formação Geral Os dados obtidos a partir das respostas à questão 9 encontram -se na Tabela 3.13 e no Gráfico 3.10.
Nessa questão – de melhor desempenho dentre as duas de Formação Geral – os \textbf{Ingressantes} tiveram média 41,0, e os Concluintes, 44,3, resultando com diferença de 3,3 pontos entre os agrupamentos.
Houve semelhança na distribuição de notas dos dois grupos, observando -se percentuais aproximados de Concluintes ( 20,1 \% ) e \textbf{Ingressantes} ( 22,5 \% ) que deixaram a questão em branco.
Ainda que não tivessem deixado a prova em branco, obtiveram nota zero na questão 1,7 \% dos \textbf{Ingressantes} e 1,6 \% dos Concluintes.
A nota mediana foi 50,0 para ambos os grupos. \textbf{Ingressantes} e Concluintes tiveram maior frequência de notas no intervalo de 60 a 70, com respectivamente 13,7 e 15,8 \% dos estudantes de cada grupo nessa faixa.
A nota máxima, 100,0 pontos, foi alcançada tanto pelos Concluintes quanto pelos \textbf{Ingressantes}. 3.3.1.2 Comentários sobre as respostas à Questão Discursiva 9 No que tange mais precisamente à atualidade do tema proposto e ao desenvolvimento do conteúdo, deve ser ressaltado que se mostrou sintonizado com a realidade política brasileira, bem como com a necessária reflexão sobre o papel na sociedade do profissional formado pelas Instituições de Ensino Superior do país.
Dissertar sobre democracia e ética foi algo apropriado para a época em que a prova foi aplicada, pois coincidiu com o período das eleições, em que tais temas são mais discutidos na vida cotidiana.
Não se pode, entretanto, deixar de se levar em conta a tendência de os estudantes terem limitado a associação do tema “ ética ” ao assunto “ política ”, considerando -se que a ética abarca não apenas a política.
Tal associação atingiu, em alguma medida, o alcance mais preciso das respostas.
As respostas apresentadas refletiram motivação para escrever sobre o tema da questão.
Alguns estudantes chegaram a desenvolver textos eloquentes e engajados.
Em alguns casos, o conceito de democracia, por exemplo, foi bem desenvolvido ao longo do texto, simultaneamente às exemplificações pedidas.
Além disso, muitos estudantes associaram o exercício da ética ao seu futuro ofício, reconhecendo o seu papel como cidadão no desempenho de sua profissão.
Essa associação insinuou, também, o esforço de reconhecimento da noção de ética prática, ainda que de maneira codificada.
Um dos pontos positivos revelou -se no fato de, na maioria das respostas analisadas, haver uma preocupação em definir a sociedade democrática.
Isso facilitou a correção e elevou a média de notas.
Outro aspecto positivo foi a presença pertinente de exposições coerentes de comportamento ético de um futuro profissional.
A articulação dos itens “ a ” e “ b ” da questão produziu, portanto, todo tipo de reflexão acerca de uma postura ética na sociedade, inclusive, corporativa.
A comprovação de que o tema ética tem sido abordado nos cursos superiores foi outro destaque positivo detectado durante as correções.
Assim, além do caráter crítico de que muitas das respostas corretas se revestiram, os relatos constantes sobre desdobramentos da ética comprovaram a abordagem sobre o tema, no mínimo, a ética profissional, nas carreiras universitárias que participaram do ENADE.
A despeito de alguns estudantes se posicionarem contra o exame, seja anulando a questão ou escrevendo algum comentário desabonador sobre o processo, foi observado que significativa parte dos participantes tem consciência político-acadêmica e percebe por que está realizando aquele tipo de prova.
Desse modo, a maioria dos estudantes tinha algo a “ dizer “ sobre o tema da questão, isto é, valia -se da oportunidade para se posicionar a respeito do assunto trazido à baila.
Ainda do ponto de vista do conteúdo desenvolvido nas respostas, observou -se preocupação solidária e consciência cidadã dos estudantes.
Mais um destaque positivo crucial se refere à habilidade textual.
Os estudantes, por exemplo, expressaram atitudes éticas e não éticas e as contrapunham ao conceito de democracia.
Há casos em que os alunos definiam a corrupção política como um exemplo de atitude não ética, enfatizando a real missão desse cargo em uma sociedade democrática : “ atender aos anseios da população e zelar pelo bem comum ”, escreveram vários estudantes.
No tocante à produção de texto dissertativo, à objetividade e ao caráter sintético, muitos estudantes corresponderam ao solicitado, elaborando pequenos textos dissertativos.
Nessas respostas, percebeu -se demonstração de domínio do assunto focado e utilização dos textos de referência como base para a elaboração.
O enunciado da questão 9 determinou que fosse elaborado um texto dissertativo que contivesse a )..., b )... e c )..., fornecendo um encaminhamento para as respostas que conduziu tanto a um resultado positivo quanto a um negativo, a saber : foi produtiva a resposta dos estudantes no sentido de que conseguiram, na maior parte das vezes, atender plenamente aos três quesitos apresentados na questão ; por outro lado, mais especificamente devido ao uso das três letras indicativas ( a, b e c ), houve tendência à elaboração de respostas divididas em itens, em vez de um texto eminentemente dissertativo.
Muitos estudantes apresentaram respostas nesse formato fragmentado para satisfazer a formalidade do roteiro sugerido no comando da questão.
A ideia-chave que deveria permear toda a resposta – democracia – se perdia em fragmentos desconexos entre si.
Se, por um lado, o roteiro do comando da questão facilitou o raciocínio de alguns dos estudantes, por outro, induziu grande parte deles a dois caminhos : i ) o das respostas curtas e objetivas, do tipo “ um profissional de saúde não ético é aquele que conta para as outras pessoas o que foi feito na consulta com o seu paciente ” e ii ) o outro caminho, que é do contraponto entre o conceito de sociedade democrática e a realidade brasileira.
Esse segundo caminho levou vários universitários a produzirem textos bem escritos, distantes, entretanto, do tema exigido.
No que se reporta às especificidades do desenvolvimento do conteúdo, a dificuldade ou falta de competência comunicativa revelada em grande parte das respostas denuncia a pouca exposição dos alunos à leitura de textos escritos de mais excelência.
Tal dificuldade mostrou -se mais evidente na impropriedade no uso de conectores e na incapacidade de elaboração de texto dissertativo, no tratamento de questões mais abstratas.
Muitos estudantes demonstraram não saber em que consistia a ética e confundiram exemplificação com definição : ao tentarem dar exemplos, via de regra, elaboravam uma definição incipiente do que se pedia.
Portanto, a dificuldade dos estudantes em argumentar e articular as ideias com coesão e coerência ainda são os grandes “ vilões ” desse processo avaliativo.
O excesso de fragmentos e textos prévios acerca do tema permitiu que a elaboração das respostas se resumisse a uma colagem ou paráfrase mal feita dos materiais disponíveis.
Em vez de subsidiar as respostas, esse material acabou por influenciar negativamente o desempenho dos estudantes.
Desse modo, verificou -se um nível de letramento ainda significativamente baixo, visto que a maioria apresentou uma visão bastante reducionista e simplificada sobre o tema em pauta.
De um modo geral, ao conceituarem ética e sociedade democrática, os estudantes demonstraram desconhecer o limite entre os dois conceitos, sem aludirem à interseção de ambos.
Foi expressiva a quantidade de redações nas quais questionavam a democracia pelo simples fato de estarem fazendo, a fortiori, aquela prova especificamente.
Em muitas das respostas, os estudantes se diziam incomodados pelo fato de estarem obrigados a realizar o exame, demonstrando sua insatisfação.
Por diversas vezes, os estudantes não foram capazes de sintetizar e de selecionar seu discurso o suficiente para atender ao que era pedido, o que ocasionou a falta de, pelo menos, um dos aspectos da resposta.
Foi constatada grande parte de respostas com problemas graves em termos textuais e gramaticais, em relação ao que se esperaria de pessoas que estão vinculadas ao Ensino ( Técnico ) Superior.
Alguns estudantes sequer conseguiam elaborar uma frase coerente e coesa, apresentando respostas que poderiam ser classificadas como de analfabetismo funcional.
De maneira ampla, manifestou -se da parte dos estudantes uma grande dificuldade no trato da língua escrita, dificuldade essa que se traduz, sobretudo, por um desrespeito bastante frequente às regras relativas : i ) à ortografia ( por exemplo : “ proficionais ”, “ infligir as leis ”, “ corrupitos ”, “ democrácia ”, “ opitar ”, “ manter o cizilo proficionau ” ), troca de “ sociedade democrática ” por “ sociedade demográfica ” e de “ ética ” por “ étnica ” ; ii ) à pontuação ; iii ) à concordância verbal ( “ os político deve (... ) eles pode ” ) ; iv ) ao modo de organização textual ( sem distinção entre o que seja dissertativo, argumentativo ou narrativo ).
Frisa -se que esse comentário não se restringe ao que se costuma chamar de “ língua padrão ”, mas ao que poderia ser chamado de “ excelência textual ”, que não se mede exclusiva ou enfaticamente pela correção gramatical, mas pela competência, objetividade, clareza e, em alguma medida, pela criatividade no ato de redigir.
Por fim, destaca -se um dos problemas dignos de nota na formulação da questão.
A transcrição do Dicionário Houaiss apresentava defeitos.
Por exemplo, a definição 5 de “ democracia ” indicava que esse é o tipo de “ governo que acata a vontade da maioria, embora respeitando os direitos e a livre expressão das minorias ”.
O emprego de um valor concessivo na oração adverbial, por oposição a um valor aditivo, contradiz as concepções efetivamente democráticas apresentadas de 1 a 4.
Repercussões desse 3.3.1.3 Análise da Questão Discursiva 10 do Componente de Formação Geral O desempenho dos estudantes na questão 10 foi inferior ao obtido na questão de número 9, com nota média 35,7 para \textbf{Ingressantes} e 38,5 para Concluintes.
A nota máxima também foi atingida em ambos os grupos, sendo que os percentuais de alunos que deixaram a questão em branco, tanto no grupo dos \textbf{Ingressantes} quanto no grupo dos Concluintes, foram ligeiramente superiores àqueles apresentados na questão 9, 26,2 \% para \textbf{Ingressantes} e 24,8 \% para Concluintes.
Houve maior percentual de alunos, considerando -se a distribuição de notas, no intervalo entre 40 e 50 pontos, tanto para \textbf{Ingressantes} ( 15,3 \% ) quanto para Concluintes ( 16,9 \% ).
Não se manifestou uma superioridade de Concluintes nas diferentes faixas de notas posteriores a esse intervalo, com equilíbrio no desempenho dos dois grupos, só ligeiramente superior no intervalo de notas entre 70 e 80 pontos.
Essas informações estão ilustradas na Tabela 3.14 e no Gráfico 3.11. valor concessivo, mal empregado, foram encontradas em incontáveis respostas, as quais diziam ser típico de um governo democrático, seguir a vontade da maioria, “ apesar ” de respeitar as minorias.
Outro desses problemas dignos de nota refere -se à indicação do dicionário quanto a “ governo ”, ao passo que o comando focava a ideia mais ampla de “ sociedade democrática ”. 3.3.1.3 Análise da Questão Discursiva 10 do Componente de Formação Geral O desempenho dos estudantes na questão 10 foi inferior ao obtido na questão de número 9, com nota média 35,7 para \textbf{Ingressantes} e 38,5 para Concluintes.
A nota máxima também foi atingida em ambos os grupos, sendo que os percentuais de alunos que deixaram a questão em branco, tanto no grupo dos \textbf{Ingressantes} quanto no grupo dos Concluintes, foram ligeiramente superiores àqueles apresentados na questão 9, 26,2 \% para \textbf{Ingressantes} e 24,8 \% para Concluintes.
Houve maior percentual de alunos, considerando -se a distribuição de notas, no intervalo entre 40 e 50 pontos, tanto para \textbf{Ingressantes} ( 15,3 \% ) quanto para Concluintes ( 16,9 \% ).
Não se manifestou uma superioridade de Concluintes nas diferentes faixas de notas posteriores a esse intervalo, com equilíbrio no desempenho dos dois grupos, só ligeiramente superior no intervalo de notas entre 70 e 80 pontos.
Essas informações estão ilustradas na Tabela 3.14 e no Gráfico 3.11. 3.3.1.4 Comentários sobre as respostas à Questão Discursiva 10 Inicialmente, será apresentada uma panorâmica sobre as respostas, com destaque para os aspectos positivos observados.
Em relação ao desenvolvimento do tema e à vinculação do mesmo à vida cotidiana dos estudantes que participaram do exame, a visão positiva em relação à Educação de Jovens e Adultos ( EJA ) reforça a ideia de que a educação pode trazer crescimento não somente profissional, mas, sobretudo, para ampliar os horizontes da cidadania.
Foram frequentes as respostas que apontavam a educação como a melhor saída, senão a única, para o resgate da cidadania.
A atuação dos governos, em especial o governo federal, foi um ponto destacado nas respostas, as quais sinalizavam para a necessidade de as instâncias governamentais assumirem um papel cada vez mais ativo na consolidação da educação – pública, sobretudo – no país.
As políticas públicas focadas na educação foram, então, um item bastante frequente nas respostas, seja para enaltecer, criticar ou requisitar a ação governamental.
Outro ponto positivo destacado diz respeito à consciência ou ao reconhecimento da precariedade de condições de trabalho dos professores, em geral, no Brasil e, em particular, daqueles ocupados com turmas de jovens e adultos.
Nesse sentido, foram inúmeras as propostas apresentadas na segunda parte das respostas com relação à imperiosidade de melhoria salarial para os docentes, e à melhoria das condições gerais de trabalho desses profissionais.
A associação direta entre valorização do professor e sucesso escolar dos estudantes foi um ponto complementar a essa preocupação com os docentes.
Os estudantes demonstraram uma visão de vínculo no processo de ensino-aprendizagem, no qual alunos e professores são partes inalienáveis entre si.
A questão da repetência e da evasão escolar, referenciadas no suporte da questão, foram então, abordadas pelos alunos por meio do viés dessa valorização docente.
Ainda que os temas específicos da repetência e da evasão não fossem contemplados a contento na maior parte das respostas, ao menos, nessa associação com a qualificação e a valorização dos professores, algo foi trabalhado analiticamente pelos estudantes.
Num plano mais abrangente, houve reconhecimento de que, no processo ensino-aprendizagem, não há atores isolados : alunos, profissionais de educação, governos e a própria sociedade como corresponsáveis pelo sucesso ou fracasso de tal processo.
A atualidade do tema e a sua explícita vinculação com o cotidiano de muitos estudantes que realizavam o exame fizeram com que os respondentes demonstrassem um visível grau de envolvimento com o tema proposto.
Ou seja, foi possível detectar um empenho notório dos estudantes ao redigir suas respostas com seriedade e comprometimento relativo à delicada problemática dos milhões de estudantes trabalhadores do país.
Utilizar todo o espaço destinado à resposta, procurando aproveitar ao máximo a oportunidade de emitir um argumento ou uma opinião revelou -se como uma das evidências do aludido empenho.
A EJA também foi reconhecida como um dos caminhos concretos para o acesso à Universidade, no médio prazo.
As respostas se voltaram, muitas vezes, à condição de depoimentos pessoais sobre a conciliação de trabalho e estudo.
Nessa direção, ao assumirem o papel de depoentes, os estudantes expuseram suas histórias de vida, suas expectativas e percalços, suas esperanças e entraves cotidianos, seus anseios mais ambiciosos e problemas estruturais, que poderiam se resolver por meio da educação.
Nesse inciso, ficava clara a relação que se formava entre educação e mobilidade social, mediatizada por iniciativas governamentais como a EJA.
Uma sugestão plausível e válida para a solução dos problemas com a EJA no país reportou -se à necessidade de maior participação do público alvo na definição de estratégias para essa modalidade de ensino.
Nesse detalhe, era notória a associação das preocupações relativas à questão 9, que trata do problema da democracia e da cidadania, com o problema específico da educação escolar.
Houve um número significativo de sugestões ou propostas para se combater o problema da repetência e da evasão, destacando -se melhores salários para os professores e aulas mais dinâmicas.
A lista dessas sugestões, contudo, era bem mais ampla e diversificada, incluindo ainda : a redução da jornada de trabalho dos jovens e adultos que participam da EJA ; melhoria de infraestrutura nas instituições responsáveis pela EJA ; incentivo financeiro aos estudantes trabalhadores por meio de um sistema de bolsas de estudo ; realização de campanhas motivacionais para o retorno aos estudos ; estímulo à participação das empresas, por meio da solicitação de recursos materiais, espaço físico para as salas, auxílio financeiro ; incremento dos cursos profissionalizantes no país ; e ampliação da oferta de ensino a distância em todo o território nacional.
Muitos estudantes abordaram, com propriedade, a associação entre um mercado de trabalho cada vez mais competitivo e a necessidade de qualificação profissional, inclusive com uma clara percepção de que os salários são tanto maiores quanto maior a quantidade de anos de estudo e/ou a qualidade da formação.
Nesse tópico, retorna -se ao papel crucial da EJA como viabilizadora dessa formação.
Em continuidade à análise, serão apresentados aspectos negativos que foram notados com frequência nas respostas à questão 10.
De um ponto de vista mais global, a dificuldade repousa no entendimento da questão, isto é, na decifração do comando.
Muitos fugiram completamente do assunto proposto, ainda que o enunciado da questão fosse farto em dados específicos ao tema.
Em vez de responderem diretamente quais eram os problemas enfrentados por aqueles que têm de trabalhar e estudar ao mesmo tempo, muitos estudantes optavam por falar das dificuldades que existem na educação brasileira, de forma estendida e pasteurizada.
Vários estudantes não conseguiam, em suas respostas, relacionar as dificuldades de conciliação da escola com o trabalho, a não ser, como já mencionado, nos casos de depoimentos pessoais, relatos de experiência, ou seja, traços de uma biografia.
Assim, eram respostas cujo teor era apenas o de registrar o abandono dos estudos por parte de vários jovens que têm de trabalhar.
O tom das respostas, nesses casos, era meramente de uma descrição e não de uma análise, muito menos de um elemento que compusesse uma argumentação.
A dificuldade de compreensão da questão se expressava nas respostas que, em lugar de mostrarem os problemas de conciliação entre escola e trabalho, prendiam -se à ideia de “ associação ” – termo que consta do enunciado da questão – entre escola e trabalho.
Muitos alunos relatavam as consequências positivas de uma boa educação no momento em que se busca bons postos de trabalho.
Tratava -se de uma fuga parcial ao comando da questão, evidenciando má compreensão do mesmo.
Uma dificuldade facilmente percebida nas respostas refere -se à tendência de os estudantes tratarem dos problemas relativos à escola e, muito separadamente, daqueles problemas relativos ao trabalho, reconhecendo uma espécie de sinergia que impediria de visualizá -los em sua individualidade, ao mesmo tempo que em sua complementaridade.
Isso significa que os estudantes distinguiam por mero isolamento, ou seja, eles não conseguiam distinguir e articular e, sobretudo, reconhecer as \textbf{convocações} mútuas entre os problemas da escola e do trabalho, no caso explícito dos estudantes trabalhadores.
Essa dificuldade de identificar, associar e analisar aspectos distintos, até mesmo concorrentes, antagônicos e complementares, denuncia a fragilidade de um pensamento simplificador que a qualidade das respostas revelou ser ainda hegemônico nas universidades brasileiras.
Notou -se a presença de preconceitos na análise dos fatores que levam ao abandono dos estudos.
Muitas respostas apontavam como causa desse abandono e da repetência a preguiça, a falta de perseverança e o cansaço dos estudantes.
Nesse tópico específico, fica evidente a responsabilização do indivíduo pelos percalços e fracassos de sua vida, descontextualizando -o de uma sociedade que, em seu movimento de conjunto, lhe impõe condicionantes.
A precariedade do uso da língua pátria de modo culto foi uma constante, conforme mencionado na análise da questão 9.
As impropriedades lexicais e a falta de elementos coesores que articulassem as partes constitutivas das respostas foram um incômodo constante.
Nesse sentido, agrava -se a situação da baixa qualidade textual em decorrência complementar da construção de textos meramente descritivos e não argumentativos, como é exigido no suporte da questão.
Isso demonstra a dificuldade dos estudantes universitários na compreensão de que uma argumentação requer, imperiosamente, justificações para os apontamentos apresentados.
Os textos das respostas eram, em sua maioria, opinativos, sem qualquer compromisso com elementos explicativos ou com a busca de justificativas. 3.3.2 Componente de Conhecimento Específico Na parte da prova relativa ao componente de Conhecimento Específico, observa -se que a diferença de desempenho foi mais acentuada entre \textbf{Ingressantes} e Concluintes.
A Tabela 3.15 contém dados elucidativos, na qual é possível observar que no grupo dos \textbf{Ingressantes} a nota média 6,7 é inferior à alcançada por Concluintes, 17,3.
Essas médias foram menores do que as obtidas no conjunto das questões objetivas, nas quais \textbf{Ingressantes} alcançaram nota média 30,2, e Concluintes, 47,3 ( Tabela 3.8 ).
O Gráfico 3.12 apresenta a distribuição de notas nos diferentes intervalos.
Observa -se que tanto \textbf{Ingressantes} ( 76,6 \% ) quanto Concluintes ( 43,5 \% ) tiveram notas, com maior frequência, no intervalo de zero a 10 pontos.
Acima desse patamar, Concluintes e \textbf{Ingressantes} apresentam proporções decrescentes com a faixa de notas.
A análise de cada uma destas questões será feita a seguir. ∗ ∗ A análise completa, com apresentação de trechos de respostas de estudantes, pode ser encontrada no documento “ Relatório da Correção das Questões Discursivas ”, elaborado pelos docentes que participaram do processo de correção das mesmas. 3.3.2.1 Análise da Questão Discursiva 38 do Componente de Conhecimento Específico Na questão 38, cujos resultados aferidos encontram -se descritas na Tabela 3.16, o grupo dos \textbf{Ingressantes} obteve nota média 2,4, enquanto a média dos Concluintes foi 9,3.
Deixaram a questão em branco 65,6 \% dos \textbf{Ingressantes} e 35,2 \% dos Concluintes.
Entre os que não deixaram a prova em branco, estes últimos obtiveram mais notas zero nesta questão do que \textbf{Ingressantes}, 32,2 \% e 24,0 \%, respectivamente.
As proporções de \textbf{Ingressantes} e Concluintes nas faixas são, grosso modo, decrescentes com a nota.
Em ambos os grupos, a nota máxima, 100,0 pontos, foi alcançada.
Entre os que não deixaram a prova em 3.3.2.2 Comentários sobre as respostas à Questão Discursiva 38 O número expressivo de questões em branco reflete a sua inadequação didático-pedagógica, fato comprovado também pelo maior índice de notas zero.
A inclusão do gráfico no enunciado da questão foi um fator que, aparentemente, confundiu os estudantes, tendo sido apontado pelos corretores como o principal responsável pelo elevado número de notas zero, visto que a tendência dos graduandos era explicar o gráfico com base no uso da eletroterapia ( ou vice-versa ).
Outra resposta muito constante foi a discussão entre amplitude e duração do pulso, com base nos valores apresentados no gráfico.
Poucos alunos responderam à questão conforme o padrão de resposta esperado, recebendo nota máxima.
Outro conjunto um pouco maior obteve notas que variaram entre 70 \% e 90 \% de aproveitamento da prova.
Esses números colocaram em evidência, de acordo com o parecer dos corretores, que a questão não teve uma formulação adequada.
Segundo os professores, a inclusão do gráfico foi desnecessária.
A solicitação da questão tinha um grau de complexidade que dificultou a elaboração de respostas em apenas 15 linhas.
Entretanto, o padrão de resposta foi elaborado com uma quantidade menor do que a solicitada para o respondente.
As respostas parcialmente corretas ou incompletas receberam notas em torno de 20 \% a 40 \% de aproveitamento da prova, mesmo sendo formuladas de maneira extensa, uma vez que, na realidade, somente pequenos trechos se referiam à pergunta.
Todos os trechos corretos das respostas foram considerados e pontuados.
Provavelmente pelo fato de tratar -se de aspectos específicos da profissão analisada, os estudantes se reportaram ao uso da eletroterapia como forma de tratamento ( equipamentos utilizados, formas de analgesia, entre outros ).
Na tentativa de escrever dentro dos parâmetros fixados, verificou -se, ainda, que os estudantes divagavam sobre o tema e só atendiam a uma parte da pergunta, o que resultou em um expressivo número de questões com notas intermediárias. 3.3.2.3 Análise da Questão Discursiva 39 do Componente de Conhecimento Específico A Tabela 3.17 contém as informações relativas à questão 39 do conjunto de questões do componente de Conhecimento Específico.
O desempenho dos dois grupos – \textbf{Ingressantes} e Concluintes – nessa questão foi relativamente superior ao desempenho na questão 38, especialmente para Concluintes.
As médias foram 3,1 no grupo dos \textbf{Ingressantes} e 21,1 no dos Concluintes.
A proporção de estudantes que deixaram a questão em branco foi semelhante à que foi observada na questão 38, sendo esses percentuais de 68,9 \% entre os \textbf{Ingressantes} e de 28,1 \% dentre os Concluintes.
A nota máxima ( 100,0 ) foi atingida pelos dois grupos.
Entre os que não deixaram a prova em branco, a nota zero foi atribuída a 19,2 \% de \textbf{Ingressantes} e a 17,7 \% de Concluintes.
Para todas as faixas correspondentes a notas não nulas, Concluintes apresentaram maiores freqüências do que os \textbf{Ingressantes}.
As distribuições das notas entre os estudantes constam do Gráfico 3.14. branco, a nota zero foi atribuída a 19,2 \% de \textbf{Ingressantes} e a 17,7 \% de Concluintes.
Para todas as faixas correspondentes a notas não nulas, Concluintes apresentaram maiores freqüências do que os \textbf{Ingressantes}.
As distribuições das notas entre os estudantes constam do Gráfico 3.14. 3.3.2.4 Comentários sobre as respostas à Questão Discursiva 39 Ficou evidente o desconhecimento dos estudantes em relação ao assunto, visto que grande parte se restringiu a divagar sobre as figuras apresentadas para conseguir responder à questão em aproximadamente 15 linhas.
A questão 39 também apresentou um número significativo de questões em branco e foi a segunda colocada em abstenção.
Apesar do desconhecimento teórico adequado sobre a utilização dos mapas, a prática clínica permitiu, em especial aos Concluintes, responder parcialmente à questão.
Dessa forma, o número de notas completamente incorretas foi bem menor em relação à questão 38.
As respostas totalmente incorretas relacionaram -se mais à tentativa de escrever a quantidade de linhas solicitada.
Apesar do número de respostas corretas, que receberam nota máxima, ter sido maior do que na questão 38, ainda assim o número foi surpreendentemente baixo.
O número de notas entre 70,0 \% e 90,0 \% de aproveitamento da prova, consideradas notas coerentes, caso a questão fosse tecnicamente mais bem elaborada, foi encontrado num número reduzido de provas. 3.3.2.5 Análise da Questão Discursiva 40 do Componente de Conhecimento Específico A Tabela 3.18 contém as informações relativas à questão 40 do conjunto do componente de Conhecimento Específico.
O desempenho dos dois agrupamentos de estudantes, nessa questão, foi superior ao ser comparado às questões 38 e 39.
A nota média de \textbf{Ingressantes} foi 14,6 ; já a de Concluintes, 21,4.
A nota mediana total foi 10,0.
Por outro lado, a nota máxima ( 100,0 ) foi obtida nos dois grupos.
A questão foi deixada em branco por 41,9 \% dos \textbf{Ingressantes} e por 24,2 \% dos Concluintes.
O intervalo de notas no qual foi observada maior frequência corresponde à faixa que vai de 10 a 20 pontos, tanto para \textbf{Ingressantes} ( com 15,5 \% ) quanto para Concluintes ( 19,8 \% ).
A distribuição de notas nos dois grupos pode ser observada no Gráfico 3.15. 3.3.2.6 Comentários sobre as respostas à Questão Discursiva 40 A correção da questão 40 evidenciou que o enunciado deu margem a múltiplos entendimentos por parte dos respondentes.
Na maioria das respostas, os estudantes manifestaram incompreensão quanto ao objetivo da pergunta e sobre quais seriam os temas a serem abordados.
Esse fato também gerou uma certa dificuldade para se constituir um processo padronizado de correção.
A questão 40, por ser extremamente abrangente quanto ao conteúdo, possibilitou uma redução na quantidade de questões em branco em relação às demais questões.
Apesar do decréscimo verificado, foi possível registrar um número expressivo de provas em branco.
O texto apresentado se relaciona à Política Nacional de Atenção Integral à Saúde do Homem e ao SUS.
A questão foi elaborada com base na saúde do homem na sociedade moderna, dividida em dois subitens extremamente amplos.
Essas incongruências metodológicas permitiram que os alunos redigissem textos calcados em qualquer dos itens abordados, tanto no texto apresentado quanto nos subitens solicitados, sem qualquer nexo.
O número de respostas totalmente corretas ou que atendessem o mais próximo possível ao padrão de resposta foi ínfimo.
A questão suscitou um elevado número de respostas parcialmente corretas em função de seu caráter abrangente.
Apesar do padrão de resposta ampliar as possibilidades de acerto, possibilitando uma correção mais justa, os estudantes consideraram apenas um ou outro item em suas respostas, o que determinou uma redução nas notas atribuídas.

% matched lemmas: convocação, ingressante
\end{textsample}
